\chapter{Hasil dan Pembahasan}

\section{Gambaran Umum Data}
Bagian ini menyajikan gambaran umum data yang digunakan.

\section{Contoh Penyajian Tabel}
Tabel~\ref{tab:miskin} menyajikan jumlah dan persentase penduduk miskin menurut
daerah. Judul tabel diletakkan di ATAS, diberi baris nomor kolom (1), (2), (3),
dan baris Sumber di bawah (untuk data sekunder).

\begin{table}[H]
\centering
% Judul tabel: DI ATAS, DIPUSATKAN (center).
\caption{Jumlah dan persentase penduduk miskin menurut daerah}
\label{tab:miskin}
% Bungkus tabel agar lebarnya diketahui -> Sumber bisa diletakkan tepat di tepi
% kiri tabel.
\sbox\skripsiblokbox{\begin{tabular}{|l|r|r|}
\hline
\textbf{Daerah} & \textbf{Jumlah (ribu orang)} & \textbf{Persentase (\%)} \\ \hline
\multicolumn{1}{|c|}{(1)} & \multicolumn{1}{c|}{(2)} & \multicolumn{1}{c|}{(3)} \\ \hline
Perkotaan  & 1.641,18 & 7,00 \\
Perdesaan  & 2.234,70 & 12,86 \\ \hline
Total      & 3.875,88 & 9,50 \\ \hline
\end{tabular}}%
\usebox\skripsiblokbox\par
\addvspace{0.2\baselineskip}
% Sumber: tepat di tepi kiri tabel.
\parbox[t]{\wd\skripsiblokbox}{\raggedright\setlength{\parindent}{0pt}\small Sumber: Badan Pusat Statistik (data ilustrasi)}
\end{table}

\section{Contoh Tabel Panjang (Lintas-Halaman)}
Bila tabel melampaui satu halaman, gunakan \texttt{xltabular} agar judul kolom dan
baris nomor kolom terulang otomatis pada setiap halaman dengan keterangan
``Tabel~N (lanjutan)''. Tabel panjang dibuat berlebar minimal 80\% lebar teks dan
dipusatkan (kolom label melebar, kolom angka tetap ringkas).

\begin{singlespace}
\setlength{\LTcapwidth}{\textwidth}   % judul DIPUSATKAN di lebar teks (center)
\begin{xltabular}{0.8\textwidth}{|r|>{\raggedright\arraybackslash}X|r|r|}
\caption{Contoh tabel panjang lintas-halaman (data ilustrasi)}\label{tab:panjang}\\
\hline
\textbf{No.} & \textbf{Kabupaten/Kota} & \textbf{Jumlah (ribu)} & \textbf{Persentase (\%)} \\ \hline
\multicolumn{1}{|c|}{(1)} & \multicolumn{1}{c|}{(2)} & \multicolumn{1}{c|}{(3)} & \multicolumn{1}{c|}{(4)} \\ \hline
\endfirsthead
\multicolumn{4}{l}{\small Tabel \thetable\ (lanjutan)}\\ \hline
\textbf{No.} & \textbf{Kabupaten/Kota} & \textbf{Jumlah (ribu)} & \textbf{Persentase (\%)} \\ \hline
\multicolumn{1}{|c|}{(1)} & \multicolumn{1}{c|}{(2)} & \multicolumn{1}{c|}{(3)} & \multicolumn{1}{c|}{(4)} \\ \hline
\endhead
\hline \endfoot
\hline \endlastfoot
1  & Kabupaten 01 & 275,7 & 15,31 \\
2  & Kabupaten 02 & 366,3 & 12,71 \\
3  & Kabupaten 03 & 536,9 & 14,74 \\
4  & Kabupaten 04 & 822,0 & 11,30 \\
5  & Kabupaten 05 & 221,4 & 10,95 \\
\end{xltabular}
\par\addvspace{0.2\baselineskip}
\centerline{\parbox[t]{0.8\textwidth}{\raggedright\setlength{\parindent}{0pt}\small Sumber: Badan Pusat Statistik (data ilustrasi)}}
\end{singlespace}

\section{Contoh Penyajian Gambar}
Gambar~\ref{gbr:bar} menyajikan persentase penduduk miskin menurut daerah. Judul
gambar diletakkan di BAWAH dan DIPUSATKAN; urutan: gambar, Sumber (tepi kiri gambar), lalu judul (pedoman).

\begin{figure}[H]
\centering
% Bungkus gambar dalam kotak agar lebarnya diketahui (untuk meletakkan Sumber
% tepat di tepi kiri gambar). Untuk gambar berkas, ganti tikzpicture dgn
% \includegraphics[...]{...}.
\sbox\skripsiblokbox{\begin{tikzpicture}
  \foreach \x/\h/\lbl in {0/2.1/Perkotaan, 2/3.9/Perdesaan, 4/2.85/Total} {
    \fill[gray!55] (\x,0) rectangle (\x+1.2,\h);
    \node[below] at (\x+0.6,0) {\small \lbl};
  }
  \draw[->] (-0.3,0) -- (5.6,0);
  \draw[->] (-0.3,0) -- (-0.3,4.4) node[above,rotate=90,anchor=south west,xshift=-3.2cm]{\small Persentase (\%)};
\end{tikzpicture}}%
\usebox\skripsiblokbox\par
\addvspace{0.15\sjudul}
% Sumber: tepat di TEPI KIRI gambar, diletakkan SEBELUM judul (pedoman).
\parbox[t]{\wd\skripsiblokbox}{\raggedright\setlength{\parindent}{0pt}\small Sumber: diolah dari Tabel~\ref{tab:miskin} (data ilustrasi)}\par
% Judul gambar: di BAWAH gambar/sumber, DIPUSATKAN (center).
\caption{Persentase penduduk miskin menurut daerah}
\label{gbr:bar}
\end{figure}

Berdasarkan Gambar~\ref{gbr:bar}, persentase kemiskinan di perdesaan lebih tinggi
dibandingkan perkotaan.

\section{Contoh Tabel Lebar (Orientasi \emph{Landscape})}
Tabel yang terlalu lebar untuk orientasi potret dapat disajikan pada halaman
\emph{landscape}, atau diperkecil hingga ukuran huruf minimal 8 pt (pedoman).
Tabel~\ref{tab:lebar} menyajikan contoh tabel lebar berorientasi \emph{landscape};
judul tetap di batas kiri teks (kiri halaman \emph{landscape}) dan sumber di
bawah-kiri, sedangkan tabel diletakkan di tengah.

\begin{landscape}
\begin{table}[H]
\centering
\caption{Contoh tabel lebar berorientasi landscape dengan banyak kolom (data ilustrasi)}
\label{tab:lebar}
\sbox\skripsiblokbox{\begin{tabular}{|l|r|r|r|r|r|r|r|r|}
\hline
\textbf{Provinsi} & \textbf{2018} & \textbf{2019} & \textbf{2020} & \textbf{2021} & \textbf{2022} & \textbf{2023} & \textbf{2024} & \textbf{2025} \\ \hline
\multicolumn{1}{|c|}{(1)} & \multicolumn{1}{c|}{(2)} & \multicolumn{1}{c|}{(3)} & \multicolumn{1}{c|}{(4)} & \multicolumn{1}{c|}{(5)} & \multicolumn{1}{c|}{(6)} & \multicolumn{1}{c|}{(7)} & \multicolumn{1}{c|}{(8)} & \multicolumn{1}{c|}{(9)} \\ \hline
Aceh        & 15,68 & 15,32 & 14,99 & 15,33 & 14,75 & 14,45 & 14,23 & 13,98 \\
Sumatera Utara & 8,94 & 8,83 & 8,75 & 9,01 & 8,42 & 8,15 & 7,99 & 7,80 \\
Jawa Barat  & 7,25 & 6,91 & 7,88 & 8,40 & 7,98 & 7,62 & 7,41 & 7,20 \\
Jawa Tengah & 11,32 & 10,80 & 11,41 & 11,79 & 10,93 & 10,77 & 10,51 & 10,30 \\
DI Yogyakarta & 11,81 & 11,44 & 12,28 & 11,91 & 11,49 & 11,04 & 10,83 & 10,62 \\ \hline
\end{tabular}}%
\usebox\skripsiblokbox\par
\addvspace{0.2\baselineskip}
\parbox[t]{\wd\skripsiblokbox}{\raggedright\setlength{\parindent}{0pt}\small Sumber: Badan Pusat Statistik (data ilustrasi)}
\end{table}
\end{landscape}
